%% ma: L11-B01-0002
%% lop: 11
%% chuong: 1
%% bai: 1
%% dang: 11.1.2
%% loai: TN4
%% muc: TH
%% dap_an: "D"
%% nguon: "de-cuong-giua-ky-1-toan-11-nam-2026-2027-truong-thpt-hoang-van-thu-ha-noi-2-8.pdf (M08, Đề cương giữa kì 1 Toán 11 THPT Hoàng Văn Thụ, Hà Nội), Phần 1 Câu 8"
%% kien_thuc: "Bài 1 (giá trị lượng giác của góc lượng giác)"
%% trang_thai: chinh-thuc
%% ghi_chu: "Đã duyệt 29/10/2026."
\begin{ex}
Rút gọn biểu thức $A=\cos\left(\dfrac{\pi}{2}-x\right)+\cos(2\pi-x)+\cos(3\pi+x)$, ta được
\choice{$A=-\cos x$}{$A=-\sin x$}{$A=\cos x$}{\True $A=\sin x$}
\loigiai{$\cos\left(\dfrac\pi2-x\right)=\sin x$, $\cos(2\pi-x)=\cos x$, $\cos(3\pi+x)=\cos(\pi+x)=-\cos x$. Vậy $A=\sin x$.}
\end{ex}
